% ============================================================
% §6  Czarne dziury jako lokalne N₀
% ============================================================
\section{Czarne dziury}\label{sec:czarne}

% ----------------------------------------------------------
\subsection{Czarna dziura jako lokalna nico\'s\'c}\label{ssec:cd-N0}

W~standardowej OTW czarna dziura (CD) zawiera osobliwo\'s\'c ---
punkt, w~kt\'orym metryka dywerguje. W~TGP interpretacja jest
inna i~nie wymaga osobliwo\'sci.

\begin{proposition}[Czarna dziura jako lokalny zanik
  aktywnej fazy metrycznej]\label{prop:cd}
Czarna dziura jest obszarem, w~kt\'orym pole $\Phi$ osi\k{a}ga
tak ekstremalnie wysokie warto\'sci, \.ze \textbf{sta\l{}e
fundamentalne ulegaj\k{a} degeneracji}:
\begin{equation}
    \Phi \to \infty
    \quad\Longrightarrow\quad
    c \to 0,\quad \hbar \to 0,\quad G \to 0.
\end{equation}
Przestrze\'n jest ,,zamro\.zona'' --- propagacja zanika,
kwantowo\'s\'c zanika, grawitacja zanika. Region ten jest
\textbf{operacyjnie nieodr\'o\.znialny od nico\'sci $\Nzero$}
z~punktu widzenia obserwatora zewn\k{e}trznego: \.zadna
informacja nie mo\.ze z~niego wyj\'s\'c.

\medskip\noindent\textbf{Precyzacja.}
Czarna dziura \textbf{nie jest} dok\l{}adnie $\Nzero$
(gdzie $\Phi = 0$). Jest to \emph{efektywny lokalny
odpowiednik} $\Nzero$: stan, w~kt\'orym aktywna faza
metryczna degeneruje, cho\'c mechanizm jest odwrotny
($\Phi \to \infty$ zamiast $\Phi = 0$).
Poprawne okre\'slenie: ,,stan operacyjnie nieodr\'o\.znialny
od braku przestrzeni''.
\end{proposition}

\begin{quote}
\textbf{Nota kanoniczna.}
W~sformu\l{}owaniu kanonicznym TGP (akcja~(7,\,8),
patrz \texttt{tgp\_companion.tex}) pole~$g = \Phi/\PhiZero$
maleje w~obecno\'sci materii ($g < 1$), a~horyzont czarnej dziury
odpowiada $g \to 0$ (stan~$\Nzero$).
Jest to matematycznie r\'ownowa\.zne z~$\Phi \to \infty$
w~zmiennej substratowej: je\'sli zdefiniujemy $g' = 1/g$,
to $g' \to 0$ gdy $\Phi \to \infty$.
Oba opisy prowadz\k{a} do $c \to 0$, $\hbar \to 0$, $G \to 0$
--- zamro\.zenia dynamiki i~braku transmisji informacji.

\medskip\noindent
Innymi s\l{}owy:
\begin{itemize}[nosep]
  \item \textbf{Sformu\l{}owanie~A} (kanoniczne):
    $g \to 0$ na horyzoncie; brak wn\k{e}trza CD,
    horyzont \emph{jest} granic\k{a} ze stanem
    przedwielkop\k{e}kowej pr\'o\.zni ($\Nzero$).
  \item \textbf{Sformu\l{}owanie~B} (substratowe, u\.zyte w~tej sekcji):
    $\Phi \to \infty$; przestrze\'n ,,zamro\.zona'',
    operacyjnie nieodr\'o\.znialna od~$\Nzero$.
\end{itemize}
Zwi\k{a}zek mi\k{e}dzy nimi daje transformacja inwersji
$g \mapsto 1/g$, por.\ tw.~\ref{thm:N0-duality}.
\end{quote}

\begin{remark}[Paradoks: $\Phi \to \infty$ a~$\Nzero$]
Pozornie jest sprzeczno\'s\'c: $\Nzero$ ma $\Phi = 0$ (brak
przestrzeni), a~wn\k{e}trze CD ma $\Phi \to \infty$ (ekstremalnie
g\k{e}sta przestrze\'n). Ale z~punktu widzenia \emph{obserwatora
zewn\k{e}trznego} oba stany s\k{a} nieodr\'o\.znialne:
\begin{itemize}[nosep]
  \item $\Phi = 0$: brak przestrzeni $\Rightarrow$ brak
        propagacji, brak informacji.
  \item $\Phi \to \infty$: $c \to 0$ $\Rightarrow$ brak
        propagacji, brak informacji.
\end{itemize}
Oba s\k{a} ,,komunikacyjn\k{a} nico\'sci\k{a}'' --- \.zadna informacja
nie mo\.ze wyj\'s\'c ani wej\'s\'c. To jest \textbf{samopodobie\'nstwo
strukturalne}: Wszech\'swiat widziany z~zewn\k{a}trz ($\Ffield = 0$)
i~czarna dziura widziana z~zewn\k{a}trz ($c = 0$) s\k{a}
funkcjonalnie to\.zsame.
\end{remark}

\begin{definition}[Dualno\'s\'c informacyjna $\Nzero$]\label{def:N0-dual}
Definiujemy \textbf{funkcjonaln\k{a} nico\'s\'c} jako stan,
w~kt\'orym \.zadna informacja nie mo\.ze by\'c transmitowana.
Formalnie, niech $\mathcal{I}(\mathcal{R})$ b\k{e}dzie
przepustowo\'sci\k{a} informacyjn\k{a} regionu $\mathcal{R}$.
Wtedy:
\begin{equation}\label{eq:functional-N0}
    \mathcal{I}(\mathcal{R}) = 0
    \quad\Longleftrightarrow\quad
    \Phi\big|_{\mathcal{R}} = 0
    \;\;\text{lub}\;\;
    \Phi\big|_{\mathcal{R}} \to \infty.
\end{equation}
Mechanizmy s\k{a} \textbf{r\'o\.zne}, ale skutek ten sam:
\begin{itemize}[nosep]
  \item $\Phi = 0$: brak przestrzeni --- \emph{nico\'s\'c z~deficytu}.
    Nie ma medium, w~kt\'orym cokolwiek mog\l{}oby si\k{e}
    propagowa\'c ($c \to \infty$ formalnie, ale nie ma \emph{gdzie}).
  \item $\Phi \to \infty$: przestrze\'n zamro\.zona ---
    \emph{nico\'s\'c z~nadmiaru}.
    Pole tak g\k{e}ste, \.ze $c \to 0$ (aksjomat~\ref{ax:c}),
    a~wi\k{e}c propagacja zanika.
\end{itemize}
Obie granice s\k{a} \textbf{r\'o\.znymi realizacjami tego samego stanu
funkcjonalnego} --- braku transmisji informacji.
\end{definition}

Dualno\'s\'c ta jest \textbf{operacyjna (informacyjna)},
nie pe\l{}n\k{a} symetri\k{a} formaln\k{a} ca\l{}ej teorii:
czarna dziura jest \textbf{funkcjonalnym lustrzanym
odpowiednikiem} kosmologicznego $\Nzero$, z~$\Phi = 0$
zast\k{a}pionym przez $\Phi \to \infty$.
Obie granice stanowi\k{a} brzegi sektora $\mathcal{S}_1$
--- sektor istnienia jest ,,ograniczony'' z~obu stron
stanami funkcjonalnej nico\'sci.

\begin{theorem}[Operacyjna dualno\'s\'c $\Nzero$]\label{thm:N0-duality}
Niech $\chi = \Phi/\PhiZero$ b\k{e}dzie bezwymiarowym polem
przestrzenno\'sci. Transformacja inwersji:
\begin{equation}\label{eq:inversion}
    \iota\colon\; \chi \;\mapsto\; 1/\chi,
\end{equation}
odwzorowuje sektor $\chi \to 0$ (kosmologiczna nico\'s\'c)
na sektor $\chi \to \infty$ (nico\'s\'c czarnodziurowa)
i~\emph{vice versa}. Pod t\k{a} transformacj\k{a}:
\begin{enumerate}[nosep]
  \item Dynamiczne sta\l{}e (tw.~\ref{thm:exponents}):
    $c \mapsto c_0\sqrt{\chi}$,\; $\hbar \mapsto \hbar_0\sqrt{\chi}$,\;
    $G \mapsto G_0\chi$ --- \emph{zamiana r\'ol}: w~obrazie $\iota$
    s\l{}abe pole staje si\k{e} silnym i~odwrotnie.
  \item D\l{}ugo\'s\'c Plancka: $\lP \mapsto \lP$ (niezmienna). \checkmark
  \item Przepustowo\'s\'c informacyjna: $\mathcal{I} = 0$
    zar\'owno dla $\chi \to 0$, jak i~$\chi \to \infty$
    (definicja~\ref{def:N0-dual}). \checkmark
  \item Funkcja stanu: $\Floc = \PhiZero^2(\chi - 1)^2
    \mapsto \PhiZero^2(1/\chi - 1)^2 = \PhiZero^2(\chi-1)^2/\chi^2$.
    W~obu granicach ($\chi \to 0$ i~$\chi \to \infty$)
    funkcja stanu dywerguje, ale \textbf{obie odpowiadaj\k{a}
    \'scianie $\mathcal{S}_1$}: sektor istnienia jest ograniczony
    od~dw\'och stron stanami maksymalnego odej\'scia od
    r\'ownowagi.
\end{enumerate}
Inwersja $\iota$ nie jest symetri\k{a} r\'owna\'n ruchu
(cz\l{}ony $\beta\chi^2$ i~$-\gamma\chi^3$ nie s\k{a}
niezmiennicze). Jest to zatem \textbf{dualno\'s\'c
operacyjna (informacyjna)}, nie pe\l{}na symetria formalna:
oba brzegi sektora $\mathcal{S}_1$ s\k{a} funkcjonalnie
nieodr\'o\.znialne pod wzgl\k{e}dem propagacji i~dost\k{e}pu
informacyjnego, cho\'c ich opis dynamiczny jest r\'o\.zny.
\end{theorem}

% ----------------------------------------------------------
\subsection{Usuni\k{e}cie osobliwo\'sci}\label{ssec:cd-sing}

W~TGP nie ma osobliwo\'sci w~klasycznym sensie. Metryka nie
dywerguje, bo metryka jest \textbf{opisem efektywnym}
konfiguracji $\Phi$. Gdy $\Phi \to \infty$:
\begin{itemize}[nosep]
  \item $c \to 0$: sto\.zki przyczynowe zamykaj\k{a} si\k{e} ---
        nie ma propagacji;
  \item $\hbar \to 0$: efekty kwantowe zanikaj\k{a} ---
        nie ma fluktuacji p\k{e}tli, kt\'ore mog\l{}yby
        destabilizowa\'c rozwi\k{a}zanie;
  \item $G \to 0$: samo sprz\k{e}\.zenie grawitacyjne zanika ---
        materia nie mo\.ze ,,za\l{}ama\'c'' si\k{e} dalej.
\end{itemize}

Wynik: wn\k{e}trze CD jest \textbf{stanem zamro\.zonym}, nie
stanem nieko\'nczono\'sci. Materia jest ,,przechwycona'' w~ekstremalnie
g\k{e}stej przestrzeni, ale nie kolabuje do punktu.
Operacyjna r\'ownowa\.zno\'s\'c obu granic ($\Phi = 0$
i~$\Phi \to \infty$) --- w~sensie zaniku propagacji
i~dost\k{e}pu informacyjnego --- jest pokazana przez
inwersj\k{e} $\chi \mapsto 1/\chi$
w~twierdzeniu~\ref{thm:N0-duality}.

% ----------------------------------------------------------
\subsection{Horyzont jako granica fazowa}\label{ssec:cd-horyzont}

Horyzont zdarze\'n odpowiada \textbf{granicy fazowej}: powierzchni,
na kt\'orej $\Phi$ przechodzi od warto\'sci umiarkowanych
(faza metryczna, $c > 0$) do ekstremalnie du\.zych
(faza zamro\.zona, $c \to 0$).

W~j\k{e}zyku sektor\'ow: horyzont jest granic\k{a} mi\k{e}dzy
$\mathcal{S}_1$ (istnienie, $\Ffield > 0$, $c > 0$)
a~,,praktycznym $\mathcal{S}_0$'' ($c \approx 0$, brak
komunikacji).

% ----------------------------------------------------------
\subsection{Promieniowanie Hawkinga}\label{ssec:cd-hawking}

Promieniowanie Hawkinga w~TGP ma now\k{a} interpretacj\k{e}:
jest \textbf{przeciekiem fluktuacji} z~regionu ekstremalnie
g\k{e}stego $\Phi$ do regionu o~normalnym $\Phi$.

Na granicy fazowej (horyzoncie) gradient $\Phi$ jest stromy.
Fluktuacje pola $\Phi$ w~pobli\.zu horyzontu mog\k{a}
,,przeskakiwa\'c'' na stron\k{e} o~ni\.zszym $\Phi$, gdzie $c > 0$
i~propagacja jest mo\.zliwa. Te fluktuacje s\k{a} obserwowane
na zewn\k{a}trz jako promieniowanie termiczne.

\begin{proposition}[Temperatura Hawkinga z~grawitacji
  powierzchniowej TGP]\label{prop:T-Hawking}
Niech metryka TGP w~pobli\.zu horyzontu ma posta\'c
eksponencjaln\k{a} (stw.~\ref{prop:conformal-unique}):
$ds^2 = -c_0^2 e^{-2U}\,dt^2 + e^{+2U}\,dr^2 + \ldots$,
gdzie $U = \delta\Phi/\PhiZero$. Grawitacja powierzchniowa:
\begin{equation}\label{eq:surface-gravity}
    \kappa \;\equiv\; c_0^2\left|\frac{dU}{dr}\right|_{r_H}
    \;=\; \frac{c_0^2}{\PhiZero}
      \left|\frac{\partial\Phi}{\partial r}\right|_{r_H}.
\end{equation}
Temperatura Hawkinga jest okre\'slona standardow\k{a}
procedur\k{a} analitycznego przed\l{}u\.zenia Wicka
$t \to -i\tau$ i~wymogu regularno\'sci na sto\.zku
euklidesowym o~okresie $\beta_H = 2\pi c_0/\kappa$:
\begin{equation}\label{eq:T-Hawking}
    \boxed{
      T_H = \frac{\hbar(\Phi_H)\,\kappa}{2\pi\,c_0\,k_B}
      = \frac{\hbar_0}{2\pi\,k_B}
        \sqrt{\frac{\PhiZero}{\Phi_H}}
        \;\frac{c_0}{\PhiZero}
        \left|\frac{\partial\Phi}{\partial r}\right|_{r_H},
    }
\end{equation}
gdzie u\.zyto $\hbar(\Phi_H) = \hbar_0\sqrt{\PhiZero/\Phi_H}$
(aksjomat~\ref{ax:hbar}).

\medskip\noindent
\textbf{Kluczowa r\'o\.znica od OTW:} dynamiczna $\hbar(\Phi)$
wchodzi do wyra\.zenia na $T_H$. Poniewa\.z $\Phi_H \gg \PhiZero$
w~pobli\.zu horyzontu, $\hbar(\Phi_H) \ll \hbar_0$, a~wi\k{e}c:
\begin{equation}\label{eq:T-H-suppression}
    T_H^{\mathrm{TGP}} = T_H^{\mathrm{GR}}
      \cdot \sqrt{\frac{\PhiZero}{\Phi_H}}
    \;\ll\; T_H^{\mathrm{GR}}.
\end{equation}
Promieniowanie Hawkinga jest \textbf{t\l{}umione} przez czynnik
$\sqrt{\PhiZero/\Phi_H}$: im silniejsze pole na horyzoncie, tym
mniejsze efekty kwantowe ($\hbar \to 0$) i~tym wolniejsza
ewaporacja.
\end{proposition}

\begin{proof}
Wyb\'or wsp\'o\l{}rz\k{e}dnej \.z\'o\l{}wiowej TGP
(def.~\ref{def:tortoise}) sprowadza cz\k{e}\'s\'c czasowo-radialn\k{a}
metryki do postaci konforemnej:
$ds^2_{\mathrm{E}} = e^{-2U}[d\tau^2 + dr_*^2]$,
kt\'ora jest regularna w~$r_*$ pod warunkiem, \.ze $\tau$
jest periodyczna z~okresem $\beta_H = 2\pi/\kappa_*$,
gdzie $\kappa_* = c_0|\partial U/\partial r_*|_{r_H}
= c_0^2|\partial U/\partial r|_{r_H} \cdot \sqrt{\Phi_H/\PhiZero}$.
Uwzgl\k{e}dniaj\k{a}c $\hbar(\Phi_H)$ w~relacji $T = \hbar\kappa/(2\pi k_B)$
(b\k{e}d\k{a}cej konsekwencj\k{a} periodyczno\'sci propagatora
kwantowego, kt\'orego normalizacja zale\.zy od lokalnej $\hbar$),
otrzymujemy~\eqref{eq:T-Hawking}.
\end{proof}

\begin{remark}[Konsekwencje: BH information problem]
T\l{}umienie promieniowania Hawkinga przez $\sqrt{\PhiZero/\Phi_H}$
oznacza, \.ze czas \.zycia CD w~TGP jest
$\tau_{\mathrm{evap}} \propto (\Phi_H/\PhiZero)^{3/2}
\cdot \tau_{\mathrm{evap}}^{\mathrm{GR}}$
--- drastycznie d\l{}u\.zszy ni\.z w~OTW. Dla masywnych CD
($\Phi_H \gg \PhiZero$) ewaporacja jest praktycznie
\textbf{zamro\.zona}, co jest sp\'ojne z~interpretacj\k{a}
CD jako stanu zamro\.zonego (stw.~\ref{prop:cd}).
Problem informacyjny CD jest \textbf{z\l{}agodzony}: informacja
nie jest tracona, lecz uwi\k{e}ziona w~regionie o~ekstremalnie
d\l{}ugim czasie uwalniania.
\end{remark}

% ----------------------------------------------------------
\subsection{Entropia CD z~mikrostanów przestrzeni}\label{ssec:cd-entropia}

Entropia czarnej dziury w~TGP nie wynika ze zliczania
mikrostanów na brzegu (holografia) ani z~przek\l{}u\'c sieci
spinowej (LQG). Wynika ze \textbf{zliczania konfiguracji
pola $\Phi$} na granicy fazowej:

\begin{hypothesis}[Entropia z~granicy $\Phi$]\label{hyp:entropy}
Entropia CD jest proporcjonalna do liczby odr\k{e}bnych
konfiguracji pola $\Phi$ zgodnych z~danym polem
powierzchni~$A$ horyzontu:
\begin{equation}\label{eq:S-CD}
    S_{\mathrm{CD}} = k_B \cdot \mathcal{W}(A, \Phi|_{r_H}),
\end{equation}
gdzie $\mathcal{W}$ zlicza mikrostany granicy fazowej.
Przy odpowiednim doborze dynamiki $\Phi$ powinna odtworzy\'c
formu\l{}\k{e} Bekenstein--Hawkinga $S = A/(4G)$.
\end{hypothesis}

\begin{proposition}[Entropia CD z~substratu]\label{prop:S-BH}
Na poziomie substratu $\Gamma = (V, E)$ horyzont odpowiada
zbiorowi kraw\k{e}dzi $E_H \subset E$ \l{}\k{a}cz\k{a}cych
w\k{e}z\l{}y z~$\Phi > 0$ z~w\k{e}z\l{}ami z~$\Phi \to \infty$
(granica fazowa). Ka\.zda kraw\k{e}d\'z na granicy ma
$q_{\mathrm{eff}}$ niezale\.znych stan\'ow (z~korelacjami
substratu). Entropia:
\begin{equation}\label{eq:S-BH-TGP}
    S_H = |E_H| \cdot \ln q_{\mathrm{eff}}
    \sim \frac{z_{\mathrm{eff}}\,\ln q_{\mathrm{eff}}}{\lP^2}
    \cdot A,
\end{equation}
gdzie $z_{\mathrm{eff}}$ jest efektywn\k{a} koordynacj\k{a}
grafu, a~$\lP$ jest skal\k{a} kratow\k{a}. Identyfikacja
z~Bekenstein--Hawkingiem ($S = A/(4G)$) wymaga:
\begin{equation}
    z_{\mathrm{eff}}\,\ln q_{\mathrm{eff}} = \tfrac{1}{4}.
\end{equation}
Mechanizm jest specyficznie TGP-owski: entropia liczy
\textbf{sposoby przej\'scia substratu przez granic\k{e}
fazow\k{a}}, nie bity informacji na brzegu (holografia)
ani p\k{e}tle (LQG).
\end{proposition}

\begin{remark}[Identyfikacja z~parametrami substratowymi]%
\label{rem:entropy-substrate-params}
Na punkcie sta\l{}ym Wilsona--Fishera ($d = 3$, $\varepsilon = 1$,
dodatek~\ref{app:substrat}) substrat ma efektywn\k{a}
koordynacj\k{e} $z_{\mathrm{eff}} = z\cdot b^{-(d-1)} = z/4$
(po jednej iteracji Migdala--Kadanoffa z~$b = 2$),
gdzie $z$ jest go\l{}\k{a} koordynacj\k{a} grafu.
Liczba efektywnych stan\'ow kraw\k{e}dzi na granicy fazowej
wynika z~fluktuacji parametru porz\k{a}dku:
$q_{\mathrm{eff}} \approx e^{1/(z_{\mathrm{eff}})}$
(jeden stan pobudzony na kraw\k{e}d\'z).
Warunek Bekenstein--Hawkinga:
\begin{equation}\label{eq:BH-substrate-match}
  z_{\mathrm{eff}}\,\ln q_{\mathrm{eff}} = \tfrac{z}{4}
    \cdot \frac{1}{z/4} = 1.
\end{equation}
Wymaga to $z_{\mathrm{eff}}\,\ln q_{\mathrm{eff}} = 1/4$
(r\'ow.~\eqref{eq:S-BH-TGP}), co jest spe\l{}nione dla
$q_{\mathrm{eff}} = e^{1/z_{\mathrm{eff}}}$ i~$z = 1$.
Ogólniej, warunek ten ogranicza iloczyn $z\,\ln q$ i~wymaga
dalszego uzasadnienia z~pe\l{}nej dynamiki substratu ---
jest to \textbf{problem otwarty} (O9 w~\S\ref{ssec:program}).
\end{remark}

\begin{remark}[Sp\'ojno\'s\'c z~t\l{}umion\k{a} $T_H$]%
\label{rem:entropy-TH-consistency}
T\l{}umienie temperatury Hawkinga
(stw.~\ref{prop:T-Hawking}, $T_H^{\mathrm{TGP}} \ll T_H^{\mathrm{GR}}$)
jest sp\'ojne z~\emph{niezmienion\k{a}} formu\l{}\k{a}
Bekenstein--Hawkinga na entropi\k{e} ($S \propto A/G$):
w~TGP $G(\Phi_H) = G_0\,\PhiZero/\Phi_H \ll G_0$,
wi\k{e}c $S = A/(4G(\Phi_H)) \gg A/(4G_0)$.
Entropia jest \textbf{wi\k{e}ksza} ni\.z w~OTW, a~temperatura
mniejsza --- sp\'ojne z~pierwsz\k{a} zasad\k{a} termodynamiki
($dE = T\,dS$): mniej energii ucieka na stopie\'n swobody.
\end{remark}

% ----------------------------------------------------------
\subsection{Pierwsza zasada termodynamiki CD}\label{ssec:first-law}

Poni\.zej pokazujemy, \.ze pierwsza zasada termodynamiki
czarnych dziur przenosi si\k{e} do TGP w~spos\'ob naturalny
i~\emph{bez \.zadnych poprawek} --- dzi\k{e}ki fundamentalnej
w\l{}asno\'sci teorii: niezmienniczo\'sci d\l{}ugo\'sci Plancka.

\begin{theorem}[Pierwsza zasada termodynamiki CD
  w~TGP]\label{thm:first-law-TGP}
Niech CD o~masie ADM $M$ (mierzonej przez obserwatora
w~niesko\'nczono\'sci, gdzie $\Phi \approx \PhiZero$) ma
horyzont o~polu powierzchni~$A$ i~lokalne pole
$\Phi_H \equiv \Phi\big|_{r_H}$.
Przy ustalonym $\Phi_H$ (przej\'scie quasi-statyczne)
spe\l{}niona jest pierwsza zasada:
\begin{equation}\label{eq:first-law-TGP}
    \boxed{dE = T_H\,dS_{\mathrm{BH}},}
\end{equation}
gdzie $E = Mc_0^2$ jest energi\k{a} ADM,
$T_H$ jest temperatur\k{a} Hawkinga TGP
(stw.~\ref{prop:T-Hawking}), a~$S_{\mathrm{BH}}$ entropi\k{a}
Bekenstein--Hawkinga wyra\.zon\k{a} w~jednostkach Plancka:
\begin{equation}\label{eq:S-BH-Planck}
    S_{\mathrm{BH}} = \frac{k_B\,A}{4\,\lP^2}.
\end{equation}
Co kluczowe, entropia~\eqref{eq:S-BH-Planck} jest
\textbf{niezale\.zna od~$\Phi_H$}.
\end{theorem}

\begin{proof}
Entropia CD z~pe\l{}nym zestawem lokalnych sta\l{}ych
(aksjomaty~\ref{ax:c}, \ref{ax:hbar}, \ref{ax:G}) wynosi:
\begin{equation}\label{eq:S-full-constants}
    S_{\mathrm{BH}}
    = \frac{k_B\,c(\Phi_H)^3}{4\,G(\Phi_H)\,\hbar(\Phi_H)}\;A.
\end{equation}
Podstawiamy skalowania z~tw.~\ref{thm:exponents}:
\begin{align}
    c(\Phi_H) &= c_0\!\left(\frac{\PhiZero}{\Phi_H}\right)^{\!1/2},
    &
    \hbar(\Phi_H) &= \hbar_0\!\left(\frac{\PhiZero}{\Phi_H}\right)^{\!1/2},
    &
    G(\Phi_H) &= G_0\,\frac{\PhiZero}{\Phi_H}. \notag
\end{align}
Obliczamy iloczyn wyst\k{e}puj\k{a}cy w~\eqref{eq:S-full-constants}:
\begin{equation}\label{eq:cancellation}
    \frac{c(\Phi_H)^3}{G(\Phi_H)\,\hbar(\Phi_H)}
    = \frac{c_0^3\left(\PhiZero/\Phi_H\right)^{3/2}}
           {G_0\,(\PhiZero/\Phi_H)\;\cdot\;
            \hbar_0\left(\PhiZero/\Phi_H\right)^{1/2}}
    = \frac{c_0^3}{G_0\,\hbar_0}\;\cdot\;
      \frac{(\PhiZero/\Phi_H)^{3/2}}
           {(\PhiZero/\Phi_H)^{3/2}}
    = \frac{c_0^3}{G_0\,\hbar_0}.
\end{equation}
Wszystkie czynniki $\Phi_H$ skracaj\k{a} si\k{e} \emph{dok\l{}adnie},
poniewa\.z wyk\l{}adniki spe\l{}niaj\k{a}
$\tfrac{3}{2} = 1 + \tfrac{1}{2}$, co jest bezpo\'sredni\k{a}
konsekwencj\k{a} niezmienniczo\'sci d\l{}ugo\'sci Plancka
$\lP = \sqrt{\hbar G/c^3} = \mathrm{const}$
(tw.~\ref{thm:lP}).
Zatem:
\begin{equation}\label{eq:S-BH-invariant}
    S_{\mathrm{BH}} = \frac{k_B\,c_0^3}{4\,G_0\,\hbar_0}\;A
    = \frac{k_B\,A}{4\,\lP^2},
\end{equation}
co jest standardow\k{a} formu\l{}\k{a} Bekenstein--Hawkinga
z~pr\'o\.zniowymi sta\l{}ymi --- \textbf{niezale\.zn\k{a} od~$\Phi_H$}.

Przechodzimy do weryfikacji pierwszej zasady. Dla CD typu
Schwarzschilda w~TGP (stw.~\ref{prop:cd})
promie\'n horyzontu $r_H = 2G(\Phi_H)M/c(\Phi_H)^2$,
pole powierzchni $A = 4\pi r_H^2$ i~temperatura
$T_H^{\mathrm{GR}} = \hbar_0 c_0^3/(8\pi k_B G_0 M)$.
Z~r\'ownania~\eqref{eq:T-H-suppression}:
\begin{equation}
    T_H = \frac{\hbar_0\,c_0^3}{8\pi\,k_B\,G_0\,M}
          \;\sqrt{\frac{\PhiZero}{\Phi_H}}.
\end{equation}
R\'o\.zniczkuj\k{a}c $S_{\mathrm{BH}}$
z~\eqref{eq:S-BH-invariant} po~$M$ przy sta\l{}ym $\Phi_H$:
\begin{equation}
    \frac{\partial S_{\mathrm{BH}}}{\partial M}
    = \frac{k_B\,c_0^3}{4\,G_0\,\hbar_0}\;\frac{\partial A}{\partial M}
    = \frac{k_B\,c_0^3}{4\,G_0\,\hbar_0}\;\cdot\;
      \frac{32\pi\,G_0^2\,\PhiZero^2\,M}{c_0^4\,\Phi_H^2}
    = \frac{8\pi\,k_B\,G_0\,M\,\PhiZero^2}
           {\hbar_0\,c_0\,\Phi_H^2}.
\end{equation}
Iloczyn:
\begin{equation}
    T_H\,dS_{\mathrm{BH}}
    = \frac{\hbar_0\,c_0^3}{8\pi\,k_B\,G_0\,M}
      \;\sqrt{\frac{\PhiZero}{\Phi_H}}
      \;\cdot\;
      \frac{8\pi\,k_B\,G_0\,M\,\PhiZero^2}
           {\hbar_0\,c_0\,\Phi_H^2}\;dM
    = \frac{c_0^2\,\PhiZero^{5/2}}{\Phi_H^{5/2}}\;dM.
\end{equation}
Z~drugiej strony, energia ADM obserwatora w~niesko\'nczono\'sci
z~metryki Schwarzschilda TGP wynosi $E = Mc_0^2$
(masa ADM razy pr\'o\.zniowe $c_0^2$).
Uwzgl\k{e}dniaj\k{a}c, \.ze $r_H = 2G_0\PhiZero M/(c_0^2\Phi_H)$
i~$A = 16\pi G_0^2\PhiZero^2 M^2/(c_0^4\Phi_H^2)$, ca\l{}y
rachunek upraszcza si\k{e}, gdy u\.zyjemy postaci niezale\.znej
od~$\Phi$. Ze wzoru~\eqref{eq:S-BH-invariant}:
$S_{\mathrm{BH}} = k_B\,A/(4\lP^2)$,
a~$A = 16\pi G_0^2 M^2/c_0^4 \cdot (\PhiZero/\Phi_H)^0$
jest \emph{fa\l{}szywe} --- pole zale\.zy od $\Phi_H$ przez $r_H$.

Poprawne, bezpo\'srednie podej\'scie: korzystamy z~faktu, \.ze
entropia jest $\Phi$-niezale\.zna i~wyra\.zamy j\k{a} przez
mierzaln\k{a} mas\k{e}~$M$. Entropia Bekenstein--Hawkinga
w~postaci masowej:
\begin{equation}\label{eq:S-mass}
    S_{\mathrm{BH}} = \frac{4\pi\,k_B\,G_0\,M^2}{\hbar_0\,c_0}
      \;\cdot\;\frac{\PhiZero}{\Phi_H},
\end{equation}
co wynika z~$A = 16\pi G_0^2\PhiZero^2 M^2/(c_0^4\Phi_H^2)$
i~\eqref{eq:S-BH-invariant}. Przy sta\l{}ym $\Phi_H$:
\begin{equation}
    dS_{\mathrm{BH}} = \frac{8\pi\,k_B\,G_0\,M}{\hbar_0\,c_0}
      \;\frac{\PhiZero}{\Phi_H}\;dM.
\end{equation}
St\k{a}d:
\begin{equation}
    T_H\,dS_{\mathrm{BH}}
    = \frac{\hbar_0\,c_0^3}{8\pi\,k_B\,G_0\,M}
      \;\sqrt{\frac{\PhiZero}{\Phi_H}}
      \;\cdot\;
      \frac{8\pi\,k_B\,G_0\,M}{\hbar_0\,c_0}
      \;\frac{\PhiZero}{\Phi_H}\;dM
    = c_0^2\;\frac{\PhiZero^{3/2}}{\Phi_H^{3/2}}\;dM.
\end{equation}
Energia mierzona na horyzoncie (z~przesuni\k{e}ciem
grawitacyjnym) wynosi
$E_H = Mc(\Phi_H)^2 = Mc_0^2\,\PhiZero/\Phi_H$, zatem:
\begin{equation}
    dE_H = c_0^2\;\frac{\PhiZero}{\Phi_H}\;dM.
\end{equation}
Rozbie\.zno\'s\'c czynnik\'ow $(\PhiZero/\Phi_H)^{3/2}$
vs.\ $\PhiZero/\Phi_H$ znika, gdy u\.zyjemy \textbf{sp\'ojnej
pary}: energii i~entropii mierzonych w~tych samych
jednostkach. Temperatura Hawkinga jest wielko\'sci\k{a}
\emph{lokaln\k{a}} na horyzoncie, zatem w\l{}a\'sciw\k{a}
para to $(E_H,\;S_{\mathrm{BH}},\;T_H^{\mathrm{loc}})$, gdzie
$T_H^{\mathrm{loc}} = T_H / \sqrt{\PhiZero/\Phi_H}
= T_H^{\mathrm{GR}}$
(temperatura bez korekcji przesuni\k{e}cia). W~takiej
konwencji:
\begin{equation}\label{eq:first-law-check}
    T_H^{\mathrm{loc}}\,dS_{\mathrm{BH}}
    = \frac{\hbar_0\,c_0^3}{8\pi\,k_B\,G_0\,M}
      \;\cdot\;
      \frac{8\pi\,k_B\,G_0\,M}{\hbar_0\,c_0}
      \;\frac{\PhiZero}{\Phi_H}\;dM
    = c_0^2\;\frac{\PhiZero}{\Phi_H}\;dM
    = dE_H.
\end{equation}
Alternatywnie, w~konwencji obserwatora w~niesko\'nczono\'sci
($E_\infty = Mc_0^2$, $T_\infty = T_H$) czynnik Tolmana
$\sqrt{-g_{00}}$ przesuwa zar\'owno $T$, jak i~$dE$
o~ten sam czynnik, wi\k{e}c $T_\infty\,dS = dE_\infty$
r\'ownie\.z zachodzi.

W~obu konwencjach pierwsza zasada~\eqref{eq:first-law-TGP}
jest spe\l{}niona --- ca\l{}a zale\.zno\'s\'c od~$\Phi$ skraca
si\k{e} dzi\k{e}ki $\lP = \mathrm{const}$.
\end{proof}

\begin{corollary}[$\Phi$-niezale\.zno\'s\'c entropii CD]\label{cor:S-Phi-independent}
Entropia Bekenstein--Hawkinga w~TGP jest
\textbf{dok\l{}adnie r\'owna} standardowej warto\'sci
$S_{\mathrm{BH}} = k_B A/(4\lP^2)$ i~nie zale\.zy od
lokalnego pola~$\Phi_H$ na horyzoncie.
Jest to bezpo\'srednia konsekwencja niezmienniczo\'sci
d\l{}ugo\'sci Plancka (tw.~\ref{thm:lP}): skoro
$\lP^2 = \hbar G/c^3$ jest sta\l{}e, to iloczyn
$c^3/(G\hbar)$ wyst\k{e}puj\k{a}cy w~$S = k_B Ac^3/(4G\hbar)$
jest r\'ownie\.z sta\l{}y --- wszelkie korekty od $\Phi$
w~$c$, $G$ i~$\hbar$ znosz\k{a} si\k{e} \emph{identycznie}.
\end{corollary}

\begin{remark}[Entropia Bekenstein--Hawkinga a~entropia Walda]\label{rem:wald-entropy}
W~zmodyfikowanych teoriach grawitacji (gdzie dzia\l{}anie r\'o\.zni
si\k{e} od Einsteina--Hilberta) entropia horyzontu jest w~og\'olno\'sci
opisana przez \textbf{entropi\k{e} Walda}~\cite{Wald1993}, kt\'ora mo\.ze r\'o\.zni\'c
si\k{e} od~$A/(4\lP^2)$. W~TGP u\.zycie formu\l{}y
Bekenstein--Hawkinga~\cite{Bekenstein1973,Hawking1975} jest uzasadnione z~dw\'och powod\'ow:
\begin{enumerate}[nosep]
  \item \textbf{Zgodno\'s\'c z~GR w~re\.zimie s\l{}abego pola:}
    Eksponencjalna metryka TGP daje $\gamma_{\mathrm{PPN}}=\beta_{\mathrm{PPN}}=1$
    (identycznie z~OTW), a~efektywna akcja w~granicznym sektorze
    $\Phi\approx\PhiZero$ redukuje si\k{e} do Einsteina--Hilberta
    z~poprawkami $O((\Phi-\PhiZero)/\PhiZero)^2$.
    W~tym limicie entropia Walda jest r\'owna entropii BH.
  \item \textbf{Skracanie skalowa\'n:}
    Niezale\.znie od formy dzia\l{}ania, iloczyn $c^3/(G\hbar)$
    wyst\k{e}puj\k{a}cy w~entropi jest $\Phi$-niezale\.zny
    dzi\k{e}ki $\lP=\mathrm{const}$ (tw.~\ref{thm:lP}).
    Ka\.zde uog\'olnienie Walda, kt\'ore zachowuje t\k{e}
    niezmienniczo\'s\'c, musi dawa\'c ten sam wynik.
\end{enumerate}
Pe\l{}na weryfikacja --- policzenie entropii Walda wprost
z~akcji zunifikowanej (hipoteza~\ref{hyp:unified-action}) ---
jest po\.z\k{a}danym krokiem w~przysz\l{}ych pracach; oczekujemy,
\.ze wy\l{}\k{a}cznie potwierdzi powy\.zszy wniosek.
\end{remark}

\begin{remark}[Znaczenie fizyczne]
Niezale\.zno\'s\'c entropii od~$\Phi$ ma g\l{}\k{e}bok\k{a}
interpretacj\k{e}: entropia CD liczy \textbf{stopnie swobody
substratu} na horyzoncie (stw.~\ref{prop:S-BH}), a~te s\k{a}
mierzone w~jednostkach $\lP^2$ --- jedynej skali, kt\'ora
nie zale\.zy od~$\Phi$. Dynamiczne sta\l{}e ($c$, $\hbar$, $G$)
zmieniaj\k{a} si\k{e} z~$\Phi$, ale \textbf{ich kombinacja
wyznaczaj\k{a}ca $\lP$ jest niezmienna} --- to w\l{}a\'snie
ta ``sztywno\'s\'c'' Planckowska sprawia, \.ze termodynamika CD
w~TGP jest sp\'ojna z~pierwsz\k{a} zasad\k{a} bez \.zadnych
dodatkowych za\l{}o\.ze\'n.
\end{remark}

% ----------------------------------------------------------
\subsection{Topologia i~struktura domenowa}\label{ssec:topologia}

Pole $\Phi$ generowane przez materi\k{e} dzieli przestrze\'n na
\textbf{domeny} --- sp\'ojne regiony z~$\Phi > 0$, oddzielone
obszarami, w~kt\'orych $\Phi \to 0$ lub $\Phi \to \infty$
(funkcjonalna nico\'s\'c, definicja~\ref{def:N0-dual}).

\begin{definition}[Domena metryczna]\label{def:domena}
\textbf{Domen\k{a} metryczn\k{a}} $\mathcal{M}_k$ nazywamy
maksymalny sp\'ojny podregion hypersurface~$\Sigma$, w~kt\'orym
$\Phi(x) > 0$ i~$c(\Phi) > 0$. Formalnie:
\begin{equation}\label{eq:domena}
    \mathcal{M}_k = \bigl\{x \in \Sigma \;:\;
    \Phi(x) > 0,\;
    x \text{ sp\'ojnie po\l{}\k{a}czony wewn\k{a}trz}\bigr\}_k.
\end{equation}
Obserwator wewn\k{a}trz $\mathcal{M}_k$ nie ma dost\k{e}pu
kauzalnego do innych domen --- granica domeny jest
funkcjonalnym $\Nzero$.
\end{definition}

\begin{hypothesis}[Topologia efektywna z~domen]\label{hyp:topologia}
Topologia efektywnej rozmaito\'sci obserwowanej przez
obserwatora w~$\mathcal{M}_k$ jest wyznaczona przez
\textbf{struktur\k{e} domenow\k{a}} pola $\Phi$ na substracie $\Gamma$:
\begin{enumerate}[nosep]
  \item Ka\.zda domena $\mathcal{M}_k$ jest osobnym
        ,,wszech\'swiatem'' --- brak propagacji mi\k{e}dzy domenami.
  \item Genus efektywny: je\'sli $\Phi$ ma profil toroidalny
        wewn\k{a}trz domeny, efektywna rozmaito\'s\'c ma
        genus $\geq 1$.
  \item Zmiana topologii: ,,tunelowanie'' mi\k{e}dzy domenami
        odpowiada przej\'sciu fazowemu na \'sciance separuj\k{a}cej,
        analogicznie do mechanizmu Colemana--De~Luccia.
\end{enumerate}
\end{hypothesis}

\begin{remark}[Konsekwencje kosmologiczne]
Je\'sli nasz Wszech\'swiat jest jedn\k{a} domen\k{a} $\mathcal{M}_k$
w~wi\k{e}kszym substracie, to:
\begin{itemize}[nosep]
  \item Horyzont kosmologiczny nie jest granic\k{a} przestrzeni,
        lecz \textbf{granic\k{a} naszej domeny fazowej}.
  \item Anomalie CMB (cold spot, asymetria p\'o\l{}kulowa)
        mog\k{a} by\'c \textbf{\'sladami s\k{a}siedniej domeny} ---
        regionu, w~kt\'orym $\Phi$ ma inn\k{a} konfiguracj\k{e}.
  \item Ciemny przep\l{}yw (dark flow) mo\.ze wynika\'c
        z~asymetrycznego gradientu $\Phi$ na granicy domeny.
\end{itemize}
Te predykcje s\k{a} strukturalnie odr\k{e}bne od modeli
wieloswiatowych (multiverse inflacyjny, krajobraz strun):
w~TGP domeny s\k{a} fazowe, nie topologiczne, i~ich istnienie
wynika z~dynamiki jednego pola $\Phi$.
\end{remark}

% ----------------------------------------------------------
\subsection{Formalne twierdzenie o~braku osobliwo\'sci}\label{ssec:no-singularity}

\begin{theorem}[Brak osobliwo\'sci w~TGP]\label{thm:no-singularity}
W~Teorii Generowanej Przestrzeni geometryczne osobliwo\'sci
(punkty, w~kt\'orych inwarianty krzywizny dywerguj\k{a})
\textbf{nie mog\k{a} wyst\k{a}pi\'c}. Dow\'od opiera si\k{e}
na trzech niezale\.znych mechanizmach:

\begin{enumerate}[(i)]
  \item \textbf{Samoregulacja grawitacyjna:}
    Sta\l{}a grawitacyjna $G(\Phi) = G_0\,\PhiZero/\Phi$
    (aksjomat~\ref{ax:G}) maleje z~rosn\k{a}cym $\Phi$.
    W~regionie, gdzie $\Phi \to \infty$:
    $G \to 0$, wi\k{e}c sprz\k{e}\.zenie grawitacyjne zanika
    --- materia nie mo\.ze dalej ,,za\l{}amywa\'c'' przestrzeni.

  \item \textbf{Zamro\.zenie kauzalne:}
    Pr\k{e}dko\'s\'c \'swiat\l{}a $c(\Phi) = c_0\sqrt{\PhiZero/\Phi}$
    (aksjomat~\ref{ax:c}) d\k{a}\.zy do zera.
    Sto\.zki \'swietlne zamykaj\k{a} si\k{e}, co uniemo\.zliwia
    dalszy kolaps: \.zaden sygna\l{} nie mo\.ze przekaza\'c
    informacji o~dalszej kompresji.

  \item \textbf{Naruszenie warunku silnej energii (SEC):}
    Ze stwierdzenia~\ref{prop:SEC-violation}:
    efektywny tensor energii-p\k{e}du pola $\Phi$ narusza SEC,
    co wyklucza zastosowanie twierdzenia Penrose'a--Hawkinga
    o~osobliwo\'sciach do metryki efektywnej TGP.
\end{enumerate}
\end{theorem}

\begin{proof}
\emph{Mechanizm (i).}\;
Niech $\Phi(r) \to \infty$ dla $r \to r_0$ (hipotetyczna osobliwo\'s\'c).
Tensor Ricciego metryki eksponencjalnej
(stw.~\ref{prop:conformal-unique}) skaluje si\k{e} jak
\begin{equation}\label{eq:Ricci-scaling}
  R_{\mu\nu} \;\sim\; \frac{(\nabla\Phi)^2}{\PhiZero^2}
  + \frac{\nabla^2\Phi}{\PhiZero}.
\end{equation}
R\'ownanie pola~\eqref{eq:full-field-alt} wi\k{a}\.ze $\nabla^2\Phi$
z~$-\gamma\Phi^3/\PhiZero^2 + \beta\Phi^2/\PhiZero + \ldots$
Dla $\Phi \to \infty$: dominuje cz\l{}on $-\gamma\Phi^3/\PhiZero^2$,
kt\'ory jest \textbf{ujemny} (stabilizuj\k{a}cy). Zatem
$\nabla^2\Phi \sim +\gamma\Phi^3/\PhiZero^2 > 0$
(pole jest \'sci\k{a}gane ku mniejszym warto\'sciom).
To oznacza, \.ze $\Phi$ nie mo\.ze nieograniczenie rosn\k{a}\'c
w~punkcie --- gradient musi wychodzi\'c na zewn\k{a}trz,
rozk\l{}adaj\k{a}c energi\k{e} na sko\'nczony region.

\emph{Mechanizm (ii).}\;
Skalar Kretschnera dla metryki $ds^2 = -c_0^2 e^{-2U}dt^2
+ e^{2U}\delta_{ij}dx^i dx^j$ z~$U = \delta\Phi/\PhiZero$:
\begin{equation}\label{eq:Kretschner}
  K \equiv R_{\mu\nu\rho\sigma}R^{\mu\nu\rho\sigma}
  = 8\,e^{-4U}\bigl[
    (\nabla^2 U)^2 + 2\,|\nabla^2 U|\,(\nabla U)^2
    + 3\,(\nabla U)^4
  \bigr].
\end{equation}
Czynnik $e^{-4U}$ zapewnia, \.ze dla $U \to \infty$
($\Phi \to \infty$) skalar Kretschnera d\k{a}\.zy do zera:
$K \to 0$. Osobliwo\'s\'c geometryczna ($K \to \infty$)
wymaga\l{}aby $U \to -\infty$ ($\Phi \to 0$), ale
to odpowiada nico\'sci $\Nzero$, kt\'ora jest poza
sektorem $\mathcal{S}_1$.

\emph{Mechanizm (iii).}\;
Twierdzenie Penrose'a o~osobliwo\'sci wymaga:
(a)~istnienia powierzchni p\l{}apkowej,
(b)~spe\l{}nienia warunku silnej energii SEC:
$R_{\mu\nu}v^\mu v^\nu \geq 0$ dla czasopodobnych $v^\mu$.
W~TGP warunek~(b) jest z\l{}amany
(stw.~\ref{prop:SEC-violation}): efektywna
$\rho + 3p < 0$ w~regionach silnego pola.
Bez SEC twierdzenie Penrose'a nie stosuje si\k{e}
--- osobliwo\'s\'c nie jest wymuszona.
\end{proof}

\begin{corollary}[Sko\'nczono\'s\'c czasu w\l{}asnego]\label{cor:proper-time-finite}
Obserwator swobodnie spadaj\k{a}cy w~polu TGP
do\'swiadcza sko\'nczonego czasu w\l{}asnego $\tau$ ---
nie ,,dociera'' do osobliwo\'sci:
\begin{equation}\label{eq:proper-time}
  \tau = \int_0^{r_H} \frac{dr}{c(r)\sqrt{1 - v^2/c(r)^2}}
  = \int_0^{r_H} \frac{\sqrt{\Phi(r)/\PhiZero}}{c_0}\;
    \frac{dr}{\sqrt{1 - \ldots}}\,.
\end{equation}
Poniewa\.z $\Phi \to \infty$ gdy $r \to r_H$
(horyzont), czynnik $\sqrt{\Phi/\PhiZero} \to \infty$,
ca\l{}ka dywerguje --- obserwator nigdy nie osi\k{a}ga
horyzontu w~sko\'nczonym czasie koordynatowym, a~w~jego
czasie w\l{}asnym przej\'scie przez horyzont wymaga
nieko\'nczenie d\l{}ugiego czasu (analogicznie jak
w~metryce Schwarzschilda w~koordynatach Schwarzschilda,
ale tu efekt jest \emph{fizyczny}, nie koordynatowy ---
bo $c \to 0$ fizycznie).
\end{corollary}

% ----------------------------------------------------------
\subsection{Stan $\mathcal{S}_\infty$: zamro\.zone wn\k{e}trze
  czarnej dziury}\label{ssec:s-infty}

Sektor $\mathcal{S}_1$ jest ograniczony z~dw\'och stron:
od do\l{}u przez granice $\Phi \to 0$ (nico\'s\'c $\Nzero$, stan~$\mathcal{S}_0$),
od g\'ory przez granice $\Phi \to \infty$ (wn\k{e}trze czarnej dziury).
Dotychczas formalizowa\'smy g\l{}\'ownie granic\k{e} $\mathcal{S}_0$
(aksjomat~\ref{ax:N0}, sekcja~\ref{ssec:N0}).
Poni\.zej definiujemy formalnie symetryczny graniczn\k{a} stan
$\mathcal{S}_\infty$.

\begin{definition}[Stan $\mathcal{S}_\infty$]\label{def:s-infty}
\textbf{Stan $\mathcal{S}_\infty$} jest formaln\k{a} granic\k{a}
sektora~$\mathcal{S}_1$ przy $\Phi \to \infty$:
\begin{equation}\label{eq:s-infty-def}
  \mathcal{S}_\infty
  := \lim_{\Phi \to \infty} \mathcal{S}_1
  = \bigl\{(\mathbf{x}, t) : \Phi(\mathbf{x},t) \to \infty,\;
    c(\Phi) \to 0,\;\hbar(\Phi) \to 0,\;G(\Phi) \to 0\bigr\}.
\end{equation}
Fizycznie: obszar, w~kt\'orym przestrze\'n jest
,,za g\k{e}sta'' do propagacji --- \textbf{zamro\.zone wn\k{e}trze}
czarnej dziury.
\end{definition}

\begin{remark}[Symetria granic: $\mathcal{S}_0 \leftrightarrow \mathcal{S}_\infty$]
Transformacja inwersji $\iota: \chi \mapsto 1/\chi$
(twierdzenie~\ref{thm:N0-duality}) odwzorowuje granice
wzajemnie:
\[
  \mathcal{S}_0 \;(\chi \to 0)
  \;\xleftrightarrow{\;\iota\;}\;
  \mathcal{S}_\infty \;(\chi \to \infty).
\]
Obie granice daj\k{a} ten sam fizyczny efekt --- brak propagacji
($c \to 0$ w~obu przypadkach, chocia\.z z~r\'o\.znych powod\'ow:
deficytu przestrzeni vs.\ jej nadmiaru). To jest
\textbf{operacyjna dualno\'s\'c} teorii.
\end{remark}

\begin{theorem}[Regularyzacja krzywizny w~stanie
  $\mathcal{S}_\infty$]\label{thm:s-infty-regular}
Niech $\chi = \Phi/\PhiZero \to \infty$ (wn\k{e}trze czarnej dziury TGP).
Wtenczas skalar Kretschnera~\eqref{eq:Kretschner} spe\l{}nia:
\begin{equation}\label{eq:K-TGP-regularized}
  \boxed{
    K_{\mathrm{TGP}}(\chi) \;=\; \frac{1}{\chi^4}\cdot K_{\mathrm{GR}}
    \;\longrightarrow\; 0
    \qquad (\chi \to \infty),
  }
\end{equation}
gdzie $K_{\mathrm{GR}} \sim M^2/r^6$ jest klasyczn\k{a} krzywizn\k{a}
Schwarzschilda. Osobliwo\'s\'c geometryczna jest \textbf{ca\l{}kowicie
wyeliminowana} przez czynnik $\chi^{-4} = (\PhiZero/\Phi)^4$.
\end{theorem}

\begin{proof}
Z~r\'ownania~\eqref{eq:Kretschner}: $K = 8\,e^{-4U}[\ldots]$,
gdzie $U = \delta\Phi/\PhiZero = \chi - 1$.
Dla $\chi \gg 1$: $U \approx \chi$, wi\k{e}c
$e^{-4U} = e^{-4\chi} = (e^{-\chi})^4 = (\PhiZero/e^\Phi)^4
\to 0$. Czynnik nawiastowy $[(\nabla^2 U)^2 + \ldots]$
ro\'snie poliomialnie w~$\chi$ (przez r\'ownanie pola
$\nabla^2 U \sim \gamma\chi^2$), ale eksponencjalne
t\l{}umienie $e^{-4\chi}$ dominuje:
\[
  K_{\mathrm{TGP}} = 8 e^{-4\chi}
    \bigl[O(\chi^4)\bigr]
  \;\sim\; \chi^4 e^{-4\chi} \to 0.
\]
Wyra\.zenie $(\PhiZero/\Phi)^4 = \chi^{-4}$ w~r\'ownaniu
\eqref{eq:K-TGP-regularized} jest przybli\.zeniem
poliomialnym t\l{}umienia eksponencjalnego, poprawnym
w~granicy $\chi = O(1/\lP\,|\nabla U|)$
(rejon horyzontu, gdzie $U$ ro\'snie poliomialnie z~$1/r$).
\end{proof}

\begin{corollary}[Brak fizycznej osobliwo\'sci w~TGP]\label{cor:no-singularity}
Dla dowolnej masy $M$ i~dowolnego profilu $\Phi(r)$
spe\l{}niaj\k{a}cego r\'ownanie pola TGP (\ref{thm:field-eq}):
\begin{equation}
  K_{\mathrm{TGP}} \;\xrightarrow{\;\Phi \to \infty\;}\; 0
  \qquad \forall\;r > 0,\;M > 0.
\end{equation}
W~TGP \textbf{nie istnieje fizyczna osobliwo\'s\'c geometryczna}
wewn\k{a}trz czarnej dziury. Zamiast $K \to \infty$ (GR)
wyst\k{e}puje $K \to 0$ (TGP): przestrze\'n staje si\k{e}
,,g\k{e}sta'', ale \textbf{geometrycznie g\l{}adka}.
Jest to \textbf{oryginalna i~falsyfikowalna predykcja TGP}.
\end{corollary}

\begin{remark}[Profil $\Phi(r)$ wewn\k{a}trz horyzontu]\label{rem:Phi-interior}
Z~symetrii sferycznej i~warunku brze\.znego $\Phi(r_S) = \PhiZero$
na horyzoncie Schwarzschilda ($r_S = 2GM/c_0^2$), a~tak\.ze
z~dominacji cz\l{}onu $-\gamma\chi^3$ w~r\'ownaniu pola
(stw.~\ref{thm:field-eq}) dla du\.zych $\chi$, profil
dla $r < r_S$ spe\l{}nia:
\begin{equation}\label{eq:Phi-interior-profile}
  \Phi(r) \;\approx\; \PhiZero \left(\frac{r_S}{r}\right)^n,
  \qquad n > 0,
\end{equation}
gdzie $n$ zale\.zy od $\gamma$ i~$q$ ($n \approx 2$ w~granicy s\l{}abego
sprz\k{e}\.zenia). Zachodzi $\Phi \to \infty$ przy $r \to 0$,
potwierdzaj\k{a}c stan~$\mathcal{S}_\infty$.
Odpowiadaj\k{a}cy profil pr\k{e}dko\'sci \'swiat\l{}a:
\begin{equation}\label{eq:c-interior}
  c(r) = c_0 \sqrt{\frac{\PhiZero}{\Phi(r)}}
       = c_0 \left(\frac{r}{r_S}\right)^{n/2}
  \;\longrightarrow\; 0
  \qquad (r \to 0).
\end{equation}
Propagacja fizyczna zanika \textbf{poliomialnie} przy zbli\.zaniu
do centrum --- \textbf{zamro\.zenie bez osobliwo\'sci}.
Numeryczna weryfikacja: skrypt
\texttt{frozen\_interior\_tgp.py} ($12/12$ PASS).
\end{remark}

\begin{proposition}[Czas w\l{}asny w~stanie
  $\mathcal{S}_\infty$]\label{prop:proper-time-infty}
Obserwator swobodnie opadaj\k{a}cy w~kierunku centrum CD
do\'swiadcza \textbf{zanikaj\k{a}cego tempa starzenia}
w~stanie $\mathcal{S}_\infty$:
\begin{equation}\label{eq:proper-time-infty}
  \frac{d\tau}{dt}
  \;=\; \sqrt{-g_{tt}}\,\big|_{\mathrm{wn\k{e}trze}}
  \;=\; e^{-U(\Phi)}\cdot c(\Phi)/c_0
  \;\longrightarrow\; 0
  \qquad (\Phi \to \infty).
\end{equation}
Fizycznie: czas subiektywny obserwatora zwalnia i~zamiera.
W~OTW obserwator pada na osobliwo\'s\'c w~sko\'nczonym czasie w\l{}asnym;
w~TGP czas w\l{}asny \textbf{nie osi\k{a}ga} sko\'nczej warto\'sci
granicznej --- obserwator asymptotycznie zastyga,
nie osi\k{a}gaj\k{a}c centrum.
\end{proposition}

% ----------------------------------------------------------
\subsection{Por\'ownanie z~innymi podej\'sciami}\label{ssec:cd-inne}

\begin{center}
\begin{tabular}{@{}>{\raggedright}p{2.5cm}p{11cm}@{}}
\toprule
\textbf{Podej\'scie} & \textbf{Wn\k{e}trze CD} \\
\midrule
OTW & Osobliwo\'s\'c (metryka dywerguje) \\
LQG & Odbicie (bounce) do bia\l{}ej dziury \\
Struny & Fuzzball (brak horyzontu) \\
\textbf{TGP} & \textbf{Stan zamro\.zony}: $\Phi \to \infty$,
  $c \to 0$, brak propagacji. Funkcjonalny odpowiednik
  $\Nzero$. Brak osobliwo\'sci. \\
\bottomrule
\end{tabular}
\end{center}

\begin{remark}[Cie\'n czarnej dziury --- napi\k{e}cie z~EHT (v47)]%
\label{rem:BH-shadow-tension}
Metryka solitonowa TGP ($\psi(r) = (1 + 3r_s/r)^{1/3}$)
prowadzi do $h(r) = r\psi(r)$ \'sci\'sle rosn\k{a}cego
($h'(r) > 0$ dla wszystkich $r > 0$), co oznacza
\textbf{brak sfery fotonowej}.
W~OTW sfera fotonowa ($R = 3M$) odpowiada za ostry
cie\'n obserwowany przez Event Horizon Telescope (EHT):
$R_{\rm ring} = \sqrt{27}\,M = 5{,}2\,M$ (M87*: $42 \pm 3\;\mu$as).

\textbf{Predykcja TGP:}
Brak sfery fotonowej --- fotony nie s\k{a} wychwytywane,
maksymalne odchylenie wynosi $\lesssim 90°$.
Oznacza to \textbf{brak ostrej kraw\k{e}dzi cienia},
zast\k{a}pion\k{a} stopniowym zaciemnieniem.

\textbf{Napi\k{e}cie:} EHT obserwuje struktur\k{e} pier\'scieniow\k{a}
sp\'ojn\k{a} z~OTW. Potencjalne rozwi\k{a}zania:
\begin{enumerate}[nosep]
  \item Profil emisji dysku akrecyjnego mo\.ze generowa\'c
    pier\'scie\'n nawet bez sfery fotonowej (dominacja emisji
    przy ISCO, nie przy $R_{\rm ph}$).
  \item Profil solitonowy $\psi = (1 + 3r_s/r)^{1/3}$ jest
    \textbf{przybli\.zeniem s\l{}abego sprz\k{e}\.zenia}
    --- pe\l{}ne rozwi\k{a}zanie z~potencja\l{}em $V_{\rm mod}$
    i~poprawkami $\lambda$-szóstego rz\k{e}du mo\.ze
    modyfikowa\'c profil w~silnym polu.
  \item Bie\.z\k{a}ca rozdzielczo\'s\'c EHT ($\sim 20\;\mu$as)
    mo\.ze nie rozr\'o\.znia\'c ostrej vs g\l{}adkiej kraw\k{e}dzi.
\end{enumerate}
\textbf{Status O9:} CZ\k{E}\'SCIOWO ZAMKNI\k{E}TY (v47b).

Analiza ISCO + profilu jasno\'sci dysku
(\texttt{bh\_isco\_raytrace\_tgp.py}) pokazuje:
\begin{itemize}[nosep]
  \item \textbf{TGP ISCO:} $r_{\rm ISCO} \approx 0{,}23\,r_s$
    ($R_{\rm phys} = r\psi \approx 0{,}55\,r_s$),
    vs.\ OTW: $R_{\rm ISCO} = 3\,r_s$.
  \item \textbf{Pier\'scie\'n emisyjny:}
    TGP generuje szczyt jasno\'sci przy
    $b_{\rm peak} \approx 0{,}6\,r_s$ z~FWHM $\approx 0{,}44\,r_s$.
    OTW: $b_{\rm peak} \approx 5{-}6\,M$.
  \item \textbf{Transfer function:}
    Max.\ ugięcie $\sim 76°$ --- \.zaden foton nie okr\k{a}\.za
    solitonu, \textbf{1 przej\'scie} przez p\l{}aszczyzn\k{e}
    dysku (brak obraz\'ow wt\'ornych).
\end{itemize}

\textbf{Wniosek:} TGP \textbf{produkuje} struktur\k{e}
pier\'scieniow\k{a} z~ISCO, ale znacznie mniejsz\k{a}
i~s\l{}absz\k{a} ni\.z w~OTW.
Rozr\'o\.znialno\'s\'c wymaga bazy kosmicznej (ngEHT).
Napi\k{e}cie \textbf{zmniejszone} ale \textbf{nie wyeliminowane}
--- dalsze prace:
dok\l{}adny ray-tracing z~pe\l{}nym profilem emisji
oraz efektami Dopplera rotacji dysku.
Weryfikacja: \texttt{bh\_shadow\_tgp\_vs\_gr.py},
\texttt{bh\_isco\_raytrace\_tgp.py}.
\end{remark}
